@@M0@@

The homomorphism @@M1@@ is obtained from the homomorphism: @@M2@@ and coincides with the one obtained from @@M3@@; the homomorphism: @@M4@@ is obtained from the commutative square (1). Property a) is therefore evident.

The homomorphism: @@M5@@ considered in b) can be defined as follows: let @@M6@@, and let @@M7@@ be an element of @@M8@@ lifting @@M9@@ (one exists since @@M10@@). To @@M11@@ one associates the image of @@M12@@ under the composite homomorphism: @@M13@@. Furthermore, let @@M14@@ be the augmentation ideal of: @@M15@@; the commutative square (2) gives a homomorphism: @@M16@@, hence a homomorphism: @@M17@@. Now @@M18@@; consequently, it suffices to check that the homomorphisms constructed are inverse to each other, which is done without difficulty.

\sgaStableTarget{sga6.e08.def.2-3}\textbf{Definition 2.3.} \uline{Let} @@M19@@ \uline{be a smoothable morphism}. \uline{The complex} @@M20@@, \uline{defined up to canonical isomorphism in the derived category} @@M21@@, \uline{is called the} relative cotangent complex \uline{of} @@M22@@ \uline{over} @@M23@@ (\uline{or of} @@M24@@).

\sgaStableTarget{sga6.e08.prop.2-4}\textbf{Proposition 2.4.} \uline{Let} @@M25@@ \uline{be a complete-intersection morphism} (1.1) \uline{which is smoothable}. \uline{Then} @@M26@@ \uline{is perfect}. \uline{More precisely}, \uline{for every}
